\documentclass[11pt]{letter}
\usepackage[utf8]{inputenc}
\usepackage[T1]{fontenc}
\usepackage{geometry}
\geometry{margin=1in}
\usepackage{url}
\usepackage{parskip}

\signature{George Teifel, Ph.D.\\
Associate Professor\\
Department of Electrical and Computer Engineering\\
University of New Mexico\\
Albuquerque, NM 87131, USA\\
gteifel@unm.edu}

\date{March 24, 2026}

\begin{document}

\begin{letter}{%
Editor-in-Chief\\
IEEE Transactions on Computer-Aided Design of\\
Integrated Circuits and Systems (TCAD)}

\opening{Dear Editor-in-Chief,}

We are pleased to submit our manuscript entitled \textbf{``Building VHDL
Processors with Claude Code: From 16-Bit to 64-Bit FPU''} for
consideration as a journal paper in IEEE Transactions on Computer-Aided
Design of Integrated Circuits and Systems.

This manuscript is an extended version of our conference paper presented
at the ACM/IEEE Workshop on Machine Learning for CAD (MLCAD 2026).  The
journal version contains more than 30\% new content beyond the
conference paper, including the following substantial additions:

\begin{enumerate}
\item \textbf{Comprehensive Comparison with Prior Work
      (Table~III, Section~V):} A new dedicated section presents a
      detailed comparison with five prior approaches to AI-assisted HDL
      generation (DAVE, Chip-Chat, HDLGen-ChatGPT, ResBench, and this
      work) across ten quantitative metrics, accompanied by
      approximately one page of narrative analysis comparing
      methodologies, qualitative differences, and trade-offs.

\item \textbf{Timing Analysis and Per-Variant Resource Breakdown
      (Section~IV.E, Table~IV):} A new section provides critical path
      analysis, maximum frequency estimation ($F_{\max} \approx
      108.7$~MHz), Verilog vs.\ VHDL synthesis comparison, and a
      detailed per-variant resource breakdown (LUTs, flip-flops, DSP48E1,
      BRAM) for all three processor variants in both languages.

\item \textbf{Extended FPU Verification Coverage
      (Section~IV.C, Table~III):} A new subsection presents an IEEE~754
      corner-case coverage matrix documenting which special values
      (denormals, $\pm\infty$, qNaN, sNaN, $\pm 0$) are tested for
      each of the 28 FPU operations via simulation and formal
      verification, along with the test vector selection methodology
      and coverage analysis.

\item \textbf{Toward Unbounded Verification (Section~IV.F):} A new
      section discusses the limitations of bounded model checking and
      presents strategies for extending verification toward unbounded
      proofs, including k-induction, abstract interpretation,
      assume-guarantee decomposition, and IC3/PDR algorithms, with a
      concrete four-phase roadmap.

\item \textbf{Extended Related Work (Section~I):} The introduction has
      been expanded with a survey of five additional recent works in
      AI-assisted hardware design: RTLCoder (Liu et al., 2024),
      ChipNeMo (NVIDIA, 2024), VeriGen (Thakur et al., 2024), RTLLM
      (Lu et al., 2024), and ChipGPT (Chang et al., 2023).

\item \textbf{Reproducibility Section (Section~II.E):} A new subsection
      details the skill file structure (1,286 lines across four files),
      the three-phase generation workflow, build automation
      infrastructure, and open-source availability.
\end{enumerate}

The new content exceeds 30\% of the total manuscript, expanding the
conference paper from approximately 10 pages to 15 pages in IEEE
journal format.  All original content from the conference paper has been
preserved and, where appropriate, strengthened with cross-references to
the new material.

The manuscript has not been submitted elsewhere and is not under
consideration by any other journal.  All authors have approved the
manuscript and agree to its submission.

We believe this work makes a significant contribution to the field of
computer-aided design by demonstrating that domain-specific context
engineering can enable AI-assisted generation of hardware designs far
exceeding the complexity of any prior work --- including a 64-bit RISC-V
processor with a full IEEE~754 double-precision floating-point unit,
verified through both simulation-based TDD and formal methods.

We respectfully request that this manuscript be considered for
publication in IEEE TCAD.

\closing{Sincerely,}

\end{letter}
\end{document}
